\subsection{Berechnung der Ausgangsspannung bei mehreren Eingangsspannungen}
% Alt: Aufgabe 5
\Aufgabe{
    \label{Aufg5}
    Gegeben ist die abgebildete Schaltung mit einem idealen Operationsverstärker. Die Betriebsspannung beträgt
    $\mathrm{\pm15\,V}$.
    \begin{itemize}
        \item[a)] Die Widerstände sind alle gleich, es gilt $R_1 = R_2 = R_3 = R_\mathrm{F} = 10\,\mathrm{k\Omega}$. Die Spannungen sind $U_\mathrm{E1} = \mathrm{1\,V}$, $U_\mathrm{E2} = \mathrm{2\,V}$,
              $U_\mathrm{E3} = \mathrm{3\,V}$. Wie hoch ist die Spannung am Ausgang $U_\mathrm{A}$?
        \item[b)] Die Eingangsspannungen bleiben gleich groß, nur die Widerstände betragen nun $R_\mathrm{F} = \mathrm{10\,k\Omega}$, $R_\mathrm{1} = \mathrm{2\,k\Omega}$,
              $R_\mathrm{2} = \mathrm{5\,k\Omega}$, $R_\mathrm{3} = \mathrm{10\,k\Omega}$. Wie hoch ist die Spannung am Ausgang $U_\mathrm{A}$?
        \item[c)] Die Eingangsspannungen bleiben gleich groß, nur die Widerstände betragen nun $R_\mathrm{F} = \mathrm{10\,k\Omega}$, $R_\mathrm{1} = \mathrm{10\,k\Omega}$,
              $R_\mathrm{2} = \mathrm{5\,k\Omega}$, $R_\mathrm{3} = \mathrm{2\,k\Omega}$. Wie hoch ist die Spannung am Ausgang $U_\mathrm{A}$?
    \end{itemize}
    \begin{figure}[H]
        \centering
        \input{Tikz/src/FigOPVmehrereEingaenge.tex}\alttext{Das Bild zeigt ein elektronisches Schaltbild eines Operationsverstärkers mit drei Eingangsspannungen, die über jeweils einen Widerstand (\(R_1\), \(R_2\), \(R_3\)) an den invertierenden Eingang des Verstärkers angeschlossen sind. Der nicht-invertierende Eingang ist mit Masse verbunden. Ein weiterer Widerstand (\(R_\mathrm{F}\)) verbindet den Ausgang des Verstärkers mit dem invertierenden Eingang als Rückkopplung. Die Spannungen \(U_{\mathrm{E1}}\), \(U_{\mathrm{E2}}\) und \(U_{\mathrm{E3}}\) fallen jeweils von oben nach unten zu Masse ab, und die Ausgangsspannung \(U_\mathrm{A}\) fällt vom Ausgang des Operationsverstärkers ebenfalls nach unten ab.}
        \label{fig:FigOPVmehrereEingaenge}
    \end{figure}

}

\Loesung{
    Formel:
    \begin{align*}
        U_\mathrm{A} = - (U_\mathrm{E1} \cdot \frac{R_\mathrm{F}}{R_1} + U_\mathrm{{E2}} \cdot \frac{R_\mathrm{F}}{R_2} + U_\mathrm{E3} \cdot \frac{R_\mathrm{F}}{R_3})
    \end{align*}
    \begin{itemize}
        \item [ \bf a)]
              \begin{align*}
                  U_\mathrm{A} =\mathrm{- (1\,V \cdot \frac{10\,k\Omega}{10\,k\Omega} + 2\,V \cdot \frac{10\,k\Omega}{10\,k\Omega} + 3\,V \cdot \frac{10\,k\Omega}{10\,k\Omega}) = -6\,V}
              \end{align*}
        \item [ \bf b)]
              \begin{align*}
                  U_\mathrm{A} =\mathrm{ - (1\,V \cdot \frac{10\,k\Omega}{2\,k\Omega} + 2\,V \cdot \frac{10\,k\Omega}{5\,k\Omega} + 3\,V \cdot \frac{10\,k\Omega}{10\,k\Omega}) = -12\,V}
              \end{align*}
        \item [ \bf c)]
              \begin{align*}
                  U_\mathrm{A} =\mathrm{ - (1\,V \cdot \frac{10\,k\Omega}{10\,k\Omega} + 2\,V \cdot \frac{10\,k\Omega}{5\,k\Omega} + 3\,V \cdot \frac{10\,k\Omega}{2\,k\Omega}) = -20\,V}
              \end{align*}
              Anmerkung: $-20\,\mathrm{V}$ nicht möglich, da die Betriebsspannung $\mathrm{\pm 15\,V}$ beträgt. Dies führt dazu, dass der niedrigste Wert $\mathrm{-15\,V}$ sein kann.
    \end{itemize}
    Siehe: Abschnitt \ref{fig: Verschaltung Verstaerker} und Tabelle \ref{tab:Grundschaltungen}
}